\begin{lemma}
\label{lemma-integral-closure-stalks}
\begin{slogan}
An element of an algebra over a ring is integral over the ring
if and only if it is locally integral at every prime ideal of the ring.
\end{slogan}
Let $\varphi : R \to S$ be a ring map.
Let $x \in S$. The following are equivalent:
\begin{enumerate}
\item $x$ is integral over $R$,
\item for every prime ideal $\mathfrak p \subset R$ the element
$x/1\in S_{\mathfrak p}$ is integral over $R_{\mathfrak p}$, and
\item for every maximal ideal $\mathfrak m\subset R$, the element
$x/1\in S_{\mathfrak m}$ is integral over $R_{\mathfrak m}$.
\end{enumerate}
\end{lemma}

\begin{proof}
The monic equation for (1) remains a monic equation after every
localization, so (1) implies (2), and (2) implies (3).
For the reverse implication let $B\subset S$ be the integral
closure of $R$ in $S$, with its original inclusion. It is an
$R$-submodule because it is an $R$-subalgebra. Lemma
\ref{lemma-integral-closure-localize} identifies its localization
$B_{\mathfrak m}\subset S_{\mathfrak m}$ with the integral closure
of $R_{\mathfrak m}$ there. Under (3), the class $\overline x$
in the quotient $R$-module $S/B$ therefore vanishes after localization
at every maximal ideal. This uses the exact quotient comparison
$(S/B)_{\mathfrak m}\cong S_{\mathfrak m}/B_{\mathfrak m}$,
which sends $(s+B)/u$ to $s/u+B_{\mathfrak m}$; the inverse is
induced by $s/u\mapsto(s+B)/u$, and the fraction equality criterion
and quotient relation prove both maps well defined and inverse.
If $\overline x\ne0$, its annihilator in $R$ is a proper ideal
and lies in a maximal ideal $\mathfrak m$. Vanishing of
$\overline x/1$ would give $u\notin\mathfrak m$ with
$u\overline x=0$, contradicting that containment. Thus
$\overline x=0$ and $x\in B$, proving (1). If $R$ is zero,
the unital map forces $S$ zero and the conclusion holds directly.

For completeness, retain also the original finite-subalgebra
calculation. Let $S'\subset S$ be generated by $\varphi(R)$ and
the original $x$. At a prime $\mathfrak p$ under (2), write its
monic equation of positive degree with coefficients
$a_i=b_i/u_i\in R_{\mathfrak p}$,
where $b_i\in R$ and $u_i\notin\mathfrak p$ for $1\leq i\leq d$.
Take $U=\prod_{i=1}^d u_i$ and the element
$$
H=\varphi(U)x^d+
\sum_{i=1}^d\varphi\left(b_i\prod_{j\ne i}u_j\right)x^{d-i}
\quad\text{of }S.
$$
Its image in $S_{\mathfrak p}$ is zero. Hence some
$v\notin\mathfrak p$ satisfies $\varphi(v)H=0$ in $S$.
Set $f=vU\notin\mathfrak p$. In $R_f$, each $u_i$ and $v$
is a unit because each is a factor of $f$. Dividing the last
equality by $vU$ in $S_f$ recovers the complete original equation
$$
x^d+\sum_{i=1}^d\varphi_f(b_i/u_i)x^{d-i}=0,
$$
where $\varphi_f:R_f\to S_f$ is the original localized map.
Thus $S'_f$, still the localization of that specified subalgebra,
is spanned over $R_f$ by $1,x,\ldots,x^{d-1}$.
The opens $D(f)$ so obtained cover $\operatorname{Spec}(R)$;
quasi-compactness (Lemma \ref{lemma-quasi-compact}) gives finitely
many $f_1,\ldots,f_n$ whose opens cover. They generate the unit
ideal, since no prime contains all of them. Lemma \ref{lemma-cover}
then shows that the original $S'$ is finite over $R$. Finally
Lemma \ref{lemma-characterize-integral} proves $x$ integral.
The factors $u_i$, $v$, $U$ and $f$ explicitly account both for
coefficient denominators and for the kernel of localization.
\end{proof}